\documentclass[11pt, letterpaper]{article}

% --- Packages ---------------------------------------------------------------
\usepackage[utf8]{inputenc}
\usepackage[T1]{fontenc}
\usepackage{mathptmx}                    % Times New Roman
\usepackage[margin=1in]{geometry}
\usepackage{graphicx}
\usepackage{amsmath, amssymb}
\usepackage{booktabs}                    % nice tables
\usepackage{longtable}                   % tables that span pages
\usepackage{xcolor}
\usepackage{hyperref}
\usepackage{listings}                    % code blocks
\usepackage{enumitem}
\usepackage{caption}
\usepackage{fancyhdr}
\usepackage{titlesec}

% --- Code listing style ----------------------------------------------------
\definecolor{codebg}{RGB}{245,245,245}
\definecolor{codegreen}{RGB}{0,128,0}
\definecolor{codegray}{RGB}{128,128,128}
\definecolor{codeblue}{RGB}{0,0,180}
\definecolor{codered}{RGB}{180,0,0}

\lstdefinestyle{pythonstyle}{
    language=Python,
    backgroundcolor=\color{codebg},
    basicstyle=\ttfamily\small,
    keywordstyle=\color{codeblue}\bfseries,
    stringstyle=\color{codered},
    commentstyle=\color{codegreen}\itshape,
    numberstyle=\tiny\color{codegray},
    numbers=left,
    numbersep=8pt,
    frame=single,
    framerule=0.4pt,
    rulecolor=\color{codegray},
    breaklines=true,
    breakatwhitespace=false,
    showstringspaces=false,
    tabsize=4,
    captionpos=b,
    aboveskip=10pt,
    belowskip=10pt,
}
\lstset{style=pythonstyle}

% --- Header / Footer ------------------------------------------------------
\pagestyle{fancy}
\fancyhf{}
\fancyhead[L]{\small SBTTD Code Documentation}
\fancyhead[R]{\small\thepage}
\renewcommand{\headrulewidth}{0.4pt}

% --- Section formatting ---------------------------------------------------
\titleformat{\section}{\Large\bfseries}{Section \thesection:}{0.5em}{}
\titleformat{\subsection}{\large\bfseries}{\thesubsection}{0.5em}{}

% --- Document metadata ----------------------------------------------------
\title{
    \textbf{SBTTD Data Analysis Code} \\[0.3em]
    \large Developer \& User Documentation \\[0.5em]
    \normalsize Supersonic Bladeless Turbine Theoretical Development
}
\author{Hadie Sabbah}
\date{Last Updated: \today}

\begin{document}

\maketitle
\tableofcontents
\newpage

% ===========================================================================
%                           INTRODUCTION
% ===========================================================================
\section{Introduction}

This document provides comprehensive documentation for the SBTTD data
analysis codebase. It is intended to serve as both a developer reference
and a user guide so that lab members can understand and use the
post-processing pipeline without needing to read the source code directly.

\subsection{Repository Structure}

The relevant code lives under:

\begin{verbatim}
2_SBTTD Data Analysis Code/
    MachPostProcessCode.py      # Main analysis script
    utils/
        __init__.py
        experimental.py         # Experimental + CFD comparison functions
        dataload_util.py        # Tecplot data import / save / load
        parameterComputation.py # Variable extraction & Reynolds number
        computation.py          # (placeholder)
        models.py               # Theoretical models (SPT, SE, geometry)
        plotting.py             # Plotting utilities
    docs/
        sbttd_code_documentation.tex   # This document
\end{verbatim}

\subsection{Conventions Used in This Document}

\begin{itemize}[nosep]
    \item Function signatures are shown in \texttt{monospace}.
    \item Parameters marked with a default value (e.g.\ \texttt{=1.0})
          are optional.
    \item All code examples assume the working directory is
          \texttt{2\_SBTTD Data Analysis Code/}.
\end{itemize}

\newpage
% ===========================================================================
%                       utils/experimental.py
% ===========================================================================
\section{\texttt{utils/experimental.py} --- Experimental Data Processing}
\label{sec:experimental}

This module provides a complete pipeline for loading wind-tunnel
experimental CSV data, auto-detecting test windows, computing statistics,
and comparing experimental results against CFD solutions extracted from
Tecplot. All functions support automatic pressure unit conversion so that
experimental and CFD data can be brought into a common unit system before
comparison.

\subsection{Design Philosophy}

The module is designed around two high-level workflow functions
(\texttt{process\_exp\_case} and \texttt{compare\_exp\_to\_cfd}) that
chain together lower-level helpers. Each function can also be called
independently. The conversion path is:

\begin{equation}
    \text{raw value}
    \;\xrightarrow{\times\;\texttt{pressure\_scale}\;+\;\texttt{pressure\_offset}}\;
    \text{native unit}
    \;\xrightarrow{\div\;\text{factor}}\;
    \text{Pa}
    \;\xrightarrow{\times\;\text{factor}}\;
    \text{target unit}
    \label{eq:conversion}
\end{equation}

This two-hop design (always passing through Pa as the canonical
intermediate) guarantees that any source unit can be converted to any
target unit without the user needing to know pairwise conversion factors.

% --- Unit Conversion Table ---------------------------------------------
\subsection{Supported Pressure Units}
\label{subsec:units}

The module stores conversion factors in a dictionary
\texttt{PRESSURE\_UNIT\_FROM\_PA}. To convert a value from Pa to the
target unit, multiply by the corresponding factor.

\begin{table}[h!]
\centering
\caption{Supported pressure units and their conversion factors from Pa.}
\label{tab:units}
\begin{tabular}{@{}llc@{}}
\toprule
\textbf{Unit key} & \textbf{Description} & \textbf{Factor (from Pa)} \\
\midrule
\texttt{'Pa'}  & Pascal              & $1.0$ \\
\texttt{'kPa'} & Kilopascal          & $10^{-3}$ \\
\texttt{'MPa'} & Megapascal          & $10^{-6}$ \\
\texttt{'bar'} & Bar                 & $10^{-5}$ \\
\texttt{'atm'} & Standard atmosphere & $1/101325$ \\
\texttt{'psi'} & Pounds per sq.\ inch & $1/6894.757$ \\
\bottomrule
\end{tabular}
\end{table}

To add a new unit, append one entry to \texttt{PRESSURE\_UNIT\_FROM\_PA}
and one to \texttt{PRESSURE\_UNIT\_LABELS} at the top of
\texttt{experimental.py}. No other code changes are needed.

% --- _convert_pressure -------------------------------------------------
\subsection{\texttt{\_convert\_pressure(value\_pa, target\_unit)}}
\label{subsec:convert_pressure}

\textbf{Purpose:} Internal helper that converts a pressure value (scalar
or NumPy array) from Pa to the specified target unit.

\begin{table}[h!]
\centering
\caption{Parameters for \texttt{\_convert\_pressure}.}
\begin{tabular}{@{}lll@{}}
\toprule
\textbf{Parameter} & \textbf{Type} & \textbf{Description} \\
\midrule
\texttt{value\_pa}   & \texttt{float} or \texttt{np.array} & Pressure in Pa \\
\texttt{target\_unit} & \texttt{str} & Key from Table~\ref{tab:units} \\
\bottomrule
\end{tabular}
\end{table}

\textbf{Returns:} The converted value in the target unit.

\textbf{Raises:} \texttt{ValueError} if \texttt{target\_unit} is not a
recognized key.

% --- set_publication_style ---------------------------------------------
\subsection{\texttt{set\_publication\_style(dpi=300)}}
\label{subsec:set_pub_style}

\textbf{Purpose:} Configure \texttt{matplotlib} rcParams for
publication-quality figures (Times New Roman, appropriate font sizes,
line widths). Call this \emph{once} at the start of your script before
any plotting.

\begin{table}[h!]
\centering
\caption{Parameters for \texttt{set\_publication\_style}.}
\begin{tabular}{@{}llll@{}}
\toprule
\textbf{Parameter} & \textbf{Type} & \textbf{Default} & \textbf{Description} \\
\midrule
\texttt{dpi} & \texttt{int} & \texttt{300} & Resolution for display and saved figures \\
\bottomrule
\end{tabular}
\end{table}

\textbf{Returns:} None (modifies global \texttt{matplotlib} state).

\begin{lstlisting}[caption={Setting publication style.}]
from utils.experimental import set_publication_style
set_publication_style(dpi=600)  # higher res for journal submission
\end{lstlisting}

% --- exp_cleaner -------------------------------------------------------
\subsection{\texttt{exp\_cleaner(dirc)}}
\label{subsec:exp_cleaner}

\textbf{Purpose:} Read an experimental CSV file produced by the data
acquisition system and return a structured dictionary with sensor data
separated by type (pressure vs.\ temperature), keyed by channel number.

The function auto-detects sensor columns using the regex pattern
\texttt{(SN\textbackslash d+)-([PT])(\textbackslash d+)}, which matches
headers like \texttt{SN1051-P01}, \texttt{SN1051-T06}, etc. Any column
that does not match this pattern is treated as metadata (e.g.\
\texttt{Frame}, \texttt{FTime[s]}, \texttt{FTime[ns]}).

\begin{table}[h!]
\centering
\caption{Parameters for \texttt{exp\_cleaner}.}
\begin{tabular}{@{}lll@{}}
\toprule
\textbf{Parameter} & \textbf{Type} & \textbf{Description} \\
\midrule
\texttt{dirc} & \texttt{str} or \texttt{Path} & Path to the CSV file \\
\bottomrule
\end{tabular}
\end{table}

\textbf{Returns:} A dictionary with four keys:

\begin{table}[h!]
\centering
\caption{Return structure for \texttt{exp\_cleaner}.}
\begin{tabular}{@{}lp{10cm}@{}}
\toprule
\textbf{Key} & \textbf{Description} \\
\midrule
\texttt{'sensor\_map'} & \texttt{dict} --- Full column metadata: \texttt{\{col\_idx: \{'serial', 'type', 'channel', 'header'\}\}} \\
\texttt{'meta'}        & \texttt{dict} --- Non-sensor columns: \texttt{\{header\_name: np.array\}} \\
\texttt{'P'}           & \texttt{dict} --- Pressure data: \texttt{\{channel\_number: np.array\}} \\
\texttt{'T'}           & \texttt{dict} --- Temperature data: \texttt{\{channel\_number: np.array\}} \\
\bottomrule
\end{tabular}
\end{table}

\begin{lstlisting}[caption={Basic usage of \texttt{exp\_cleaner}.}]
from utils.experimental import exp_cleaner
data = exp_cleaner(r"C:\...\Scan.csv")
print(data['P'].keys())     # e.g. dict_keys([1, 2, 3, ...])
print(data['meta'].keys())  # e.g. dict_keys(['Frame', 'FTime[s]', ...])
\end{lstlisting}

% --- detect_test_window -----------------------------------------------
\subsection{\texttt{detect\_test\_window(clean\_data, k\_sigma=10, n\_baseline=50)}}
\label{subsec:detect_test_window}

\textbf{Purpose:} Automatically detect the start and end of the
wind-tunnel test run from the pressure time series. The algorithm:

\begin{enumerate}[nosep]
    \item Identifies the stagnation channel (highest mean pressure).
    \item Estimates the noise floor from the first \texttt{n\_baseline}
          frames.
    \item Sets a threshold at $\mu_{\text{baseline}} + k_\sigma \cdot
          \sigma_{\text{baseline}}$.
    \item Returns the first and last frame indices where pressure exceeds
          this threshold.
\end{enumerate}

\begin{table}[h!]
\centering
\caption{Parameters for \texttt{detect\_test\_window}.}
\begin{tabular}{@{}llll@{}}
\toprule
\textbf{Parameter} & \textbf{Type} & \textbf{Default} & \textbf{Description} \\
\midrule
\texttt{clean\_data} & \texttt{dict}  & ---        & Output from \texttt{exp\_cleaner()} \\
\texttt{k\_sigma}    & \texttt{float} & \texttt{10}  & Threshold in standard deviations \\
\texttt{n\_baseline} & \texttt{int}   & \texttt{50}  & Number of pre-test frames for noise estimation \\
\bottomrule
\end{tabular}
\end{table}

\textbf{Returns:} \texttt{(i\_start, i\_end, t\_start, t\_end)} ---
frame indices and corresponding times in seconds.

\textbf{Raises:} \texttt{ValueError} if no frames exceed the threshold.

% --- exp_plotter ------------------------------------------------------
\subsection{\texttt{exp\_plotter(x\_data, y\_dict, y\_name, label\_name, y\_unit, show=True)}}
\label{subsec:exp_plotter}

\textbf{Purpose:} Plot time-series data for each sensor channel on a
single axes. Primarily used to visualize raw pressure or temperature
traces during a test run.

\begin{table}[h!]
\centering
\caption{Parameters for \texttt{exp\_plotter}.}
\begin{tabular}{@{}llll@{}}
\toprule
\textbf{Parameter} & \textbf{Type} & \textbf{Default} & \textbf{Description} \\
\midrule
\texttt{x\_data}     & \texttt{np.array} & ---          & Time array (x-axis) \\
\texttt{y\_dict}     & \texttt{dict}     & ---          & \texttt{\{channel: np.array\}} \\
\texttt{y\_name}     & \texttt{str}      & ---          & Y-axis label name (e.g.\ ``Pressure'') \\
\texttt{label\_name} & \texttt{str}      & ---          & Legend prefix (e.g.\ ``P'' or ``T'') \\
\texttt{y\_unit}     & \texttt{str}      & ---          & Unit string (e.g.\ ``Pa'') \\
\texttt{show}        & \texttt{bool}     & \texttt{True} & Whether to call \texttt{plt.show()} \\
\bottomrule
\end{tabular}
\end{table}

\textbf{Returns:} \texttt{(fig, ax)} --- Matplotlib figure and axes
objects for further customization.

% --- exp_stats_pressure -----------------------------------------------
\subsection{\texttt{exp\_stats\_pressure(clean\_data, confidence=0.95)}}
\label{subsec:exp_stats_pressure}

\textbf{Purpose:} Compute descriptive statistics for each pressure
channel over the auto-detected test window. Internally calls
\texttt{detect\_test\_window()} to crop the data.

Statistics computed for each channel:

\begin{itemize}[nosep]
    \item $\bar{x}$ (mean)
    \item $s$ (standard deviation, with Bessel's correction $\text{ddof}=1$)
    \item $\text{SEM} = s / \sqrt{N}$
    \item 95\% confidence interval: $\text{CI} = t_{\alpha/2,\,N-1}
          \cdot \text{SEM}$
\end{itemize}

\begin{table}[h!]
\centering
\caption{Parameters for \texttt{exp\_stats\_pressure}.}
\begin{tabular}{@{}llll@{}}
\toprule
\textbf{Parameter} & \textbf{Type} & \textbf{Default} & \textbf{Description} \\
\midrule
\texttt{clean\_data} & \texttt{dict}  & ---       & Output from \texttt{exp\_cleaner()} \\
\texttt{confidence}  & \texttt{float} & \texttt{0.95} & Confidence level for CI \\
\bottomrule
\end{tabular}
\end{table}

\textbf{Returns:} \texttt{dict} ---
\texttt{\{channel: \{'mean', 'std', 'sem', 'CI\_95', 'N'\}\}}.

Note: the returned values are in the \emph{raw sensor units} (whatever
the CSV contains). Unit conversion happens downstream in
\texttt{process\_exp\_case}.

% --- slice_extractor_wall ---------------------------------------------
\subsection{\texttt{slice\_extractor\_wall(dirc)}}
\label{subsec:slice_extractor_wall}

\textbf{Purpose:} Extract the bottom-wall surface data from a 3D
Tecplot CFD solution by slicing at $z = 0$. Requires an active Tecplot
session (\texttt{tp.session.connect()}).

\textbf{Algorithm:}
\begin{enumerate}[nosep]
    \item Load the Tecplot binary file (or reuse if already loaded).
    \item Set the frame to Cartesian 3D.
    \item Extract a slice at $z = 0$ from the ``Section'' surface zone,
          using \texttt{OneZonePerConnectedRegion} mode.
    \item Select the zone with the lowest mean $Y$ (the bottom wall).
    \item Sort all arrays by $X$ and return.
\end{enumerate}

\begin{table}[h!]
\centering
\caption{Parameters for \texttt{slice\_extractor\_wall}.}
\begin{tabular}{@{}lll@{}}
\toprule
\textbf{Parameter} & \textbf{Type} & \textbf{Description} \\
\midrule
\texttt{dirc} & \texttt{str} or \texttt{Path} & Path to \texttt{mcfd\_tec.bin} \\
\bottomrule
\end{tabular}
\end{table}

\textbf{Returns:} \texttt{(x\_wall, y\_wall, p\_wall, T\_wall)} --- all
\texttt{np.array}, sorted by $x$. Pressures are in the CFD solver's
native unit (typically Pa).

% --- process_exp_case -------------------------------------------------
\subsection{\texttt{process\_exp\_case(...)}}
\label{subsec:process_exp_case}

\textbf{Purpose:} One-call pipeline that loads a CSV, computes elapsed
time, converts pressures, runs statistics, and returns everything in a
single consistent pressure unit. This is the primary entry point for
processing an experimental dataset.

\textbf{Full signature:}
\begin{lstlisting}[numbers=none]
process_exp_case(
    csv_path,
    tap_spacing_mm,
    skip_channels=0,
    pressure_scale=1.0,
    pressure_offset=0.0,
    exp_native_unit='Pa',
    pressure_unit='Pa',
    plot=True,
)
\end{lstlisting}

\begin{longtable}{@{}lllp{6.5cm}@{}}
\caption{Parameters for \texttt{process\_exp\_case}.}
\label{tab:process_exp_case_params} \\
\toprule
\textbf{Parameter} & \textbf{Type} & \textbf{Default} & \textbf{Description} \\
\midrule
\endfirsthead
\toprule
\textbf{Parameter} & \textbf{Type} & \textbf{Default} & \textbf{Description} \\
\midrule
\endhead
\texttt{csv\_path}        & \texttt{str/Path} & ---           & Path to experimental CSV \\
\texttt{tap\_spacing\_mm} & \texttt{float}    & ---           & Physical distance between taps [mm] \\
\texttt{skip\_channels}   & \texttt{int}      & \texttt{0}    & Leading channels to exclude from stats \\
\texttt{pressure\_scale}  & \texttt{float}    & \texttt{1.0}  & Multiplier: raw $\to$ \texttt{exp\_native\_unit} \\
\texttt{pressure\_offset} & \texttt{float}    & \texttt{0.0}  & Offset added after scaling (in native unit) \\
\texttt{exp\_native\_unit}& \texttt{str}      & \texttt{'Pa'} & Unit produced by scale + offset \\
\texttt{pressure\_unit}   & \texttt{str}      & \texttt{'Pa'} & Desired output unit (see Table~\ref{tab:units}) \\
\texttt{plot}             & \texttt{bool}     & \texttt{True} & Show time-series plots \\
\bottomrule
\end{longtable}

\textbf{Unit conversion logic} (see also Eq.~\ref{eq:conversion}):

\begin{enumerate}[nosep]
    \item Multiply raw sensor reading by \texttt{pressure\_scale}.
    \item Add \texttt{pressure\_offset}. The result is now in
          \texttt{exp\_native\_unit}.
    \item Convert to Pa: divide by
          \texttt{PRESSURE\_UNIT\_FROM\_PA[exp\_native\_unit]}.
    \item Convert Pa to \texttt{pressure\_unit}: multiply by
          \texttt{PRESSURE\_UNIT\_FROM\_PA[pressure\_unit]}.
\end{enumerate}

\textbf{Returns:} A dictionary with the following keys:

\begin{table}[h!]
\centering
\caption{Return structure for \texttt{process\_exp\_case}.}
\begin{tabular}{@{}lp{9cm}@{}}
\toprule
\textbf{Key} & \textbf{Description} \\
\midrule
\texttt{'clean\_data'}    & Raw output from \texttt{exp\_cleaner()} \\
\texttt{'t\_elapsed'}     & Elapsed time array [s] \\
\texttt{'P\_converted'}   & \texttt{\{channel: np.array\}} in \texttt{pressure\_unit} \\
\texttt{'T\_data'}        & \texttt{\{channel: np.array\}} temperature data \\
\texttt{'stats'}          & Raw statistics from \texttt{exp\_stats\_pressure()} \\
\texttt{'channels'}       & Sorted channel list (after skipping) \\
\texttt{'means'}          & \texttt{np.array} of mean pressures in \texttt{pressure\_unit} \\
\texttt{'errors'}         & \texttt{np.array} of 95\% CI values in \texttt{pressure\_unit} \\
\texttt{'x\_exp\_mm'}     & Tap locations in mm \\
\texttt{'pressure\_unit'} & String recording the output unit \\
\bottomrule
\end{tabular}
\end{table}

\subsubsection*{Example: Flat Plate (raw in bar, output in kPa)}

\begin{lstlisting}[caption={Processing a flat-plate experiment.}]
from utils.experimental import set_publication_style, process_exp_case

set_publication_style()

flat = process_exp_case(
    csv_path=r"C:\...\Scan.csv",
    tap_spacing_mm=7.25,
    skip_channels=3,       # skip stagnation ports
    pressure_scale=1e5,    # raw is in bar -> * 1e5 = Pa
    pressure_offset=101325,# gauge -> absolute
    exp_native_unit='Pa',  # after scale+offset we have Pa
    pressure_unit='kPa',   # output everything in kPa
)
print(flat['means'])       # array in kPa
print(flat['pressure_unit'])  # 'kPa'
\end{lstlisting}

\subsubsection*{Example: h/l = 0.02 Angled (raw in atm, output in bar)}

\begin{lstlisting}[caption={Processing an angled-geometry experiment.}]
angled = process_exp_case(
    csv_path=r"C:\...\Test1.csv",
    tap_spacing_mm=5.65,
    skip_channels=3,
    pressure_scale=101325,  # raw is in atm -> * 101325 = Pa
    exp_native_unit='Pa',
    pressure_unit='bar',    # output in bar
)
\end{lstlisting}

% --- plot_exp_errorbar ------------------------------------------------
\subsection{\texttt{plot\_exp\_errorbar(exp\_result, ...)}}
\label{subsec:plot_exp_errorbar}

\textbf{Purpose:} Generate an error-bar plot of mean wall pressure vs.\
tap location from the output of \texttt{process\_exp\_case}. The
$y$-axis label is auto-generated from the stored
\texttt{pressure\_unit} unless overridden.

\begin{table}[h!]
\centering
\caption{Parameters for \texttt{plot\_exp\_errorbar}.}
\begin{tabular}{@{}llll@{}}
\toprule
\textbf{Parameter} & \textbf{Type} & \textbf{Default} & \textbf{Description} \\
\midrule
\texttt{exp\_result} & \texttt{dict}  & ---                              & Output from \texttt{process\_exp\_case()} \\
\texttt{title}       & \texttt{str}   & \texttt{"Experimental Wall Pressure"} & Plot title \\
\texttt{xlabel}      & \texttt{str}   & \texttt{"X [mm]"}                & X-axis label \\
\texttt{ylabel}      & \texttt{str}   & \texttt{None}                    & Auto from unit if \texttt{None} \\
\texttt{show}        & \texttt{bool}  & \texttt{True}                    & Call \texttt{plt.show()} \\
\bottomrule
\end{tabular}
\end{table}

\textbf{Returns:} \texttt{(fig, ax)}.

\begin{lstlisting}[caption={Plotting experimental error bars.}]
from utils.experimental import plot_exp_errorbar

fig, ax = plot_exp_errorbar(flat, title="Flat Plate")
# ylabel will automatically be "Pressure [kPa]"
\end{lstlisting}

% --- compare_exp_to_cfd ----------------------------------------------
\subsection{\texttt{compare\_exp\_to\_cfd(...)}}
\label{subsec:compare_exp_to_cfd}

\textbf{Purpose:} Load CFD wall data from Tecplot, trim to a matching
$x$-range, convert both CFD and experimental pressures to a common unit,
and produce an overlay plot. This function guarantees unit consistency
between the two data sources.

\textbf{Full signature:}
\begin{lstlisting}[numbers=none]
compare_exp_to_cfd(
    exp_result, cfd_path, x_min_cfd, x_max_cfd,
    cfd_native_unit='Pa',
    cfd_pressure_offset=0.0,
    pressure_unit=None,
    normalize=False,
    title="Pressure Vs X",
    xlabel="X [mm]",
    ylabel=None,
    show=True,
)
\end{lstlisting}

\begin{longtable}{@{}lllp{5.8cm}@{}}
\caption{Parameters for \texttt{compare\_exp\_to\_cfd}.}
\label{tab:compare_params} \\
\toprule
\textbf{Parameter} & \textbf{Type} & \textbf{Default} & \textbf{Description} \\
\midrule
\endfirsthead
\toprule
\textbf{Parameter} & \textbf{Type} & \textbf{Default} & \textbf{Description} \\
\midrule
\endhead
\texttt{exp\_result}           & \texttt{dict}       & ---           & Output from \texttt{process\_exp\_case()} \\
\texttt{cfd\_path}             & \texttt{str/Path}   & ---           & Path to \texttt{mcfd\_tec.bin} \\
\texttt{x\_min\_cfd}           & \texttt{float}      & ---           & Left trim bound (CFD coords) \\
\texttt{x\_max\_cfd}           & \texttt{float}      & ---           & Right trim bound (CFD coords) \\
\texttt{cfd\_native\_unit}     & \texttt{str}        & \texttt{'Pa'} & Unit the CFD solver writes \\
\texttt{cfd\_pressure\_offset} & \texttt{float}      & \texttt{0.0}  & Subtracted from CFD $P$ before conversion (in native unit) \\
\texttt{pressure\_unit}        & \texttt{str/None}   & \texttt{None} & Target unit; if \texttt{None}, matches experiment \\
\texttt{normalize}             & \texttt{bool}       & \texttt{False}& Normalize both by their maxima \\
\texttt{title}                 & \texttt{str}        & ---           & Plot title \\
\texttt{xlabel}                & \texttt{str}        & ---           & X-axis label \\
\texttt{ylabel}                & \texttt{str/None}   & \texttt{None} & Auto from unit if \texttt{None} \\
\texttt{show}                  & \texttt{bool}       & \texttt{True} & Call \texttt{plt.show()} \\
\bottomrule
\end{longtable}

\textbf{Unit matching logic:}
\begin{itemize}[nosep]
    \item If \texttt{pressure\_unit=None}: the function reads
          \texttt{exp\_result['pressure\_unit']} and converts CFD data to
          that same unit. \textbf{No mismatch is possible.}
    \item If \texttt{pressure\_unit='psi'} (or any valid key): both CFD
          and experimental data are converted to that unit, regardless of
          what either started in.
    \item The CFD conversion path is:
          $P_{\text{CFD}} - \texttt{offset} \;\to\; \text{Pa} \;\to\;
          \texttt{pressure\_unit}$.
    \item If the experiment was already processed in a different unit, it
          is re-converted: $\texttt{exp\_unit} \;\to\; \text{Pa} \;\to\;
          \texttt{pressure\_unit}$.
\end{itemize}

\textbf{Returns:} A dictionary:

\begin{table}[h!]
\centering
\caption{Return structure for \texttt{compare\_exp\_to\_cfd}.}
\begin{tabular}{@{}lp{8.5cm}@{}}
\toprule
\textbf{Key} & \textbf{Description} \\
\midrule
\texttt{'x\_cfd'}, \texttt{'P\_cfd'}         & Trimmed CFD arrays in \texttt{pressure\_unit} \\
\texttt{'x\_exp'}, \texttt{'P\_exp'}, \texttt{'errors'} & Experimental arrays in \texttt{pressure\_unit} \\
\texttt{'x\_wall'}, \texttt{'y\_wall'}, \texttt{'P\_wall'}, \texttt{'T\_wall'} & Raw wall data (CFD native unit) \\
\texttt{'pressure\_unit'}                     & The unit used for comparison \\
\texttt{'fig'}, \texttt{'ax'}                 & Matplotlib figure and axes \\
\bottomrule
\end{tabular}
\end{table}

\subsubsection*{Example: Auto-match experiment unit}

\begin{lstlisting}[caption={CFD vs.\ experiment --- auto unit matching.}]
from utils.experimental import compare_exp_to_cfd

# flat was processed with pressure_unit='kPa'
result = compare_exp_to_cfd(
    flat,
    cfd_path=r"C:\...\mcfd_tec.bin",
    x_min_cfd=40, x_max_cfd=127,
    cfd_native_unit='Pa',
    cfd_pressure_offset=94748,
    # pressure_unit=None -> auto uses 'kPa' from flat
)
# Both curves are in kPa, ylabel = "P [kPa]"
\end{lstlisting}

\subsubsection*{Example: Force a different unit}

\begin{lstlisting}[caption={Override to compare in psi.}]
result = compare_exp_to_cfd(
    flat,
    cfd_path=r"C:\...\mcfd_tec.bin",
    x_min_cfd=40, x_max_cfd=127,
    cfd_native_unit='Pa',
    cfd_pressure_offset=94748,
    pressure_unit='psi',  # forces both to psi
)
# ylabel = "P [psi]"
\end{lstlisting}

\subsubsection*{Example: Normalized comparison}

\begin{lstlisting}[caption={Normalized overlay.}]
result = compare_exp_to_cfd(
    flat,
    cfd_path=r"C:\...\mcfd_tec.bin",
    x_min_cfd=40, x_max_cfd=127,
    cfd_native_unit='Pa',
    cfd_pressure_offset=94748,
    normalize=True,
)
# ylabel = "Normalized P [-]"
\end{lstlisting}

% ===========================================================================
%                           COMPLETE WORKFLOW
% ===========================================================================
\newpage
\section{Complete Workflow Example}
\label{sec:workflow}

The following script demonstrates the full pipeline for processing three
experimental cases and comparing each to its corresponding CFD solution,
all in a common pressure unit.

\begin{lstlisting}[caption={Full workflow for three experimental cases.}]
from pathlib import Path
from utils.experimental import (
    set_publication_style,
    process_exp_case,
    plot_exp_errorbar,
    compare_exp_to_cfd,
)

set_publication_style()

# -- 1. Flat Plate ------------------------------------------
flat = process_exp_case(
    csv_path=r"C:\...\Scan.csv",
    tap_spacing_mm=7.25,
    skip_channels=3,
    pressure_scale=1e5,
    pressure_offset=101325,
    exp_native_unit='Pa',
    pressure_unit='kPa',
)
plot_exp_errorbar(flat, title="Flat Plate")

compare_exp_to_cfd(
    flat,
    cfd_path=r"C:\...\mcfd_tec.bin",
    x_min_cfd=40, x_max_cfd=127,
    cfd_native_unit='Pa',
    cfd_pressure_offset=94748,
)

# -- 2. h/l = 0.02 Angled ----------------------------------
angled = process_exp_case(
    csv_path=r"C:\...\Test1.csv",
    tap_spacing_mm=5.65,
    skip_channels=3,
    pressure_scale=101325,
    exp_native_unit='Pa',
    pressure_unit='kPa',
)
plot_exp_errorbar(angled, title="h/l = 0.02 Angled")

compare_exp_to_cfd(
    angled,
    cfd_path=r"C:\...\mcfd_tec.bin",
    x_min_cfd=-0.005, x_max_cfd=0.0678,
    cfd_native_unit='Pa',
    cfd_pressure_offset=101325,
)

# -- 3. h/l = 0.02 Straight --------------------------------
straight = process_exp_case(
    csv_path=r"C:\...\P_tap_results_Test8.csv",
    tap_spacing_mm=7.5,
    skip_channels=1,
    pressure_scale=101325,
    exp_native_unit='Pa',
    pressure_unit='kPa',
)
plot_exp_errorbar(straight, title="h/l = 0.02 Straight")
\end{lstlisting}

% ===========================================================================
%                           ADDING NEW UNITS
% ===========================================================================
\section{How to Add a New Pressure Unit}
\label{sec:add_unit}

To add support for a new pressure unit (e.g.\ mmHg), edit the two
dictionaries at the top of \texttt{utils/experimental.py}:

\begin{lstlisting}[caption={Adding mmHg support.}]
PRESSURE_UNIT_FROM_PA = {
    'Pa':   1.0,
    'kPa':  1e-3,
    # ... existing entries ...
    'mmHg': 1.0 / 133.322,   # <-- add this line
}

PRESSURE_UNIT_LABELS = {
    'Pa': 'Pa',
    # ... existing entries ...
    'mmHg': 'mmHg',           # <-- add this line
}
\end{lstlisting}

No other code changes are required. The new unit will work with all
functions immediately.

\newpage
% ===================================================================
%                       GUI APPLICATION
% ===================================================================
\section{\texttt{sweep\_study\_gui.py} --- Desktop GUI Application}
\label{sec:gui}

This section documents the graphical user interface (GUI) application
that wraps the analysis pipeline into a point-and-click tool for lab
members.

\subsection{Prerequisites}

\begin{lstlisting}[numbers=none]
pip install customtkinter
\end{lstlisting}

All other dependencies (numpy, matplotlib, scipy) should already be
installed from the analysis code.

\subsection{Launching the GUI}

\begin{lstlisting}[numbers=none]
cd "2_SBTTD Data Analysis Code"
python sweep_study_gui.py
\end{lstlisting}

\subsection{GUI Architecture}

The application is built on three core concepts:

\begin{enumerate}[nosep]
    \item \textbf{Widgets}: Visual elements (buttons, text fields,
          dropdowns, labels). Each is a Python object you create and
          place on the window.
    \item \textbf{Layout (Grid)}: Widgets are placed in a grid (row,
          column). Think of the window as a spreadsheet.
    \item \textbf{Callbacks}: Functions triggered by user actions. E.g.,
          clicking ``Browse'' opens a file dialog.
\end{enumerate}

The basic pattern for every widget:
\begin{lstlisting}[numbers=none]
widget = SomeWidget(parent, options...)    # create
widget.grid(row=R, column=C)              # place
widget.configure(command=my_function)     # connect callback
\end{lstlisting}

\subsection{Application Structure}

The GUI is a single class (\texttt{SweepStudyApp}) inheriting from
\texttt{customtkinter.CTk}:

\begin{table}[h!]
\centering
\caption{GUI layout structure.}
\begin{tabular}{@{}llp{7cm}@{}}
\toprule
\textbf{Region} & \textbf{Widget} & \textbf{Purpose} \\
\midrule
Left sidebar     & \texttt{CTkFrame}    & Study type selector, theme toggle \\
Center (Tab 1)   & ``Experimental Data'' & CSV file selection, sensor parameters, pressure conversion settings \\
Center (Tab 2)   & ``CFD Comparison''    & Tecplot file, domain trim bounds, CFD unit settings \\
Center (Tab 3)   & ``Results''           & Embedded matplotlib plot area \\
Bottom           & Status bar            & Feedback messages \\
\bottomrule
\end{tabular}
\end{table}

\subsection{How to Extend the GUI}

\subsubsection*{Adding a New Input Field}

To add a new parameter (e.g., a custom title for the plot):

\begin{lstlisting}[caption={Adding a new text entry to a tab.}]
# Inside _build_exp_tab(), add these lines:
ctk.CTkLabel(tab, text="Plot Title:").grid(
    row=row, column=0, padx=10, pady=5, sticky="w")
self.plot_title_var = ctk.StringVar(value="My Plot")
ctk.CTkEntry(tab, textvariable=self.plot_title_var, width=300).grid(
    row=row, column=1, padx=5, pady=5, sticky="w")
row += 1
\end{lstlisting}

Then read the value in the callback:
\begin{lstlisting}[numbers=none]
title = self.plot_title_var.get()
\end{lstlisting}

\subsubsection*{Adding a New Tab}

\begin{lstlisting}[caption={Adding a new tab.}]
# In _create_main_area():
self.tab_new = self.tabview.add("Batch Processing")

# Then create a builder method:
def _build_batch_tab(self):
    tab = self.tab_new
    # ... add widgets here ...
\end{lstlisting}

\subsubsection*{Adding a New Button with a Callback}

\begin{lstlisting}[caption={Button with callback pattern.}]
def _build_my_section(self):
    ctk.CTkButton(
        tab,
        text="Export to CSV",
        command=self._export_results  # <-- function called on click
    ).grid(row=row, column=0, pady=10)

def _export_results(self):
    """This runs when the user clicks 'Export to CSV'."""
    # ... do work here ...
    self._set_status("Exported successfully!")
\end{lstlisting}

\subsubsection*{Embedding a Matplotlib Plot}

The key technique for scientific GUIs:

\begin{lstlisting}[caption={Embedding matplotlib inside tkinter.}]
from matplotlib.backends.backend_tkagg import FigureCanvasTkAgg

# 1. Create figure (do NOT call plt.show())
fig, ax = plt.subplots(figsize=(8, 5))
ax.plot(x, y)

# 2. Wrap in canvas and attach to a frame
canvas = FigureCanvasTkAgg(fig, master=self.plot_frame)
canvas.draw()
canvas.get_tk_widget().pack(fill="both", expand=True)
\end{lstlisting}

\subsection{User Workflow}

\begin{enumerate}
    \item Launch the app (\texttt{python sweep\_study\_gui.py}).
    \item In the \textbf{Experimental Data} tab:
    \begin{enumerate}[nosep]
        \item Click ``Browse'' and select your CSV file.
        \item Set tap spacing, skip channels, and pressure conversion
              parameters.
        \item Choose your desired output unit from the dropdown.
        \item Click ``Process Experimental Data''.
    \end{enumerate}
    \item View the error-bar plot in the \textbf{Results} tab.
    \item (Optional) In the \textbf{CFD Comparison} tab:
    \begin{enumerate}[nosep]
        \item Browse for the \texttt{mcfd\_tec.bin} file.
        \item Set the $x$-trim bounds and pressure offset.
        \item Choose a comparison unit or leave on ``Auto''.
        \item Click ``Compare Experiment vs CFD''.
    \end{enumerate}
    \item The overlay plot appears in the \textbf{Results} tab.
\end{enumerate}

\end{document}
